\begin{proof}

\begin{enumerate}[label={$\langle1\rangle$\arabic*.},leftmargin=*]
\item Suppose its Euler number were one. The rational
normal-crossing case has Euler number two by P4.8, so we
are in the simply connected resolved-conductor alternative
of P4.26. That proof gives $\pi_1(S)=\pi_1(C)=F_b$,
$b\le1$, and
$e(C)=c+1-b$ for the number $c$ of normalized rational
components. Since $e(S_{\rm reg})\ge0$, the value one
forces $c=b=1$ and $e(S_{\rm reg})=0$.
Thus $C$ has normalization $E=\mathbb P^1$, exactly one
point with two analytic branches, and only unibranch
singular points elsewhere.

\item P4.33 proves that $N$ has rational singularities and
that its image conductor scheme has
$\chi(\mathcal O_{C_{\rm sch}})=-1$.
The generic ordinary-triple conductor and pure
Cohen--Macaulay argument of P4.31 identify this scheme
with the symbolic double curve $\mathcal O_Y/P^{(2)}$.
At a point with two analytic branches, the symbolic ideal
means order at least two along each branch. The same
argument applies: a nonzero finite-support defect in the
conductor quotient would violate depth one.

\item We compute its Euler characteristic and the possible
local defects. On $E$ use the bundles $K$ and the
ramification divisor $R$ of P4.30. This time
$\deg K=2+\deg R$, since $E=\mathbb P^1$.
Write $\tau_p$ for the local length of the conormal
quotient $j_*K/(P/P^{(2)})$ and $r_p$ for the sum of
the coefficients of $R$ over $p$. Normalization and
Riemann--Roch give
\[
 \chi(\mathcal O_Y/P^{(2)})
 =5-\sum_p d_p,
 \qquad d_p=\delta(C,p)+\tau_p-r_p.
\]
We prove the precise gap needed here: at an ordinary
node $d_p=5$; at every planar singularity
$d_p=5\delta(C,p)$; and at every nonplanar singularity
having at most two branches $d_p\ge7$.

\item First consider a planar curve, with ideal $(z,h(x,y))$.
Its conormal is free over $\mathcal O_C$, generated by
$dz,dh$. On each normalized branch, write
$dh=a_i u_i$ for a primitive plane kernel generator,
and put $v_i=\operatorname{ord}a_i$.
The inclusion of each of these two free summands into
its normalized module has length $\delta(C,p)$.
The second summand has the additional vanishing orders
$v_i$. Hence
\[
 \tau_p=2\delta(C,p)+\sum_i v_i.
\]
We verify $\sum_i(v_i-r_i)=2\delta(C,p)$ directly by
curve duality. Let $A$ be the reduced plane-curve ring,
$B$ its normalization, and $I=\operatorname{Hom}_A(B,A)$
its conductor. Applying $\operatorname{Hom}_A(-,A)$ to
$0\to A\to B\to B/A\to0$ gives
$A/I\simeq\operatorname{Ext}^1_A(B/A,A)$.
Here $\operatorname{Ext}^1_A(B,A)=0$ because $B$ is
maximal Cohen--Macaulay over the one-dimensional
Gorenstein ring $A$. The dual of a finite-length module
has the same length, so $\operatorname{length}(A/I)=\delta$,
and $\operatorname{length}(B/I)=2\delta$.
If the conductor exponent on branch $i$ is $c_i$,
this says $\sum_i c_i=2\delta$.
Choose a generic plane coordinate $x$ with minimal
valuation on every normalized branch. The dualizing
generator $dx/h_y$ then pulls back with
order $r_i-v_i$. The identity
$\omega_B=I\omega_A$ therefore gives
$c_i=v_i-r_i$. Combining these equalities proves
$\tau_p-r_p=4\delta$ and the stated planar formula.

\item Now consider a nonplanar unibranch curve of multiplicity
$m$. Its embedding dimension is three. As in P4.30,
all gradient entries have valuation at least $m$, so
$\tau_p\ge2m$, while $r_p=m-1$ and
$\delta\ge m-1$. A nonplanar unibranch curve has
$m\ge3$. To see the last assertion, if $x=t^2$, write
each other coordinate as $a(x)+t b(x)$. After subtracting
the even parts, choose the nonzero $b$ of least order.
The other odd coordinate is an analytic multiple of it
over $\mathbb C\{x\}$. This is a smooth ambient equation
containing the curve, so its embedding dimension is at
most two. The case where both odd parts vanish contradicts
the birational normalization parametrization.

\item The estimates give $d_p\ge2m\ge6$. Equality at six
would require $m=3$, $\delta=2$, and $\tau_p=6$.
The valuation semigroup of a unibranch curve has exactly
$\delta$ gaps, as follows by filtering its local ring
and its normalization by order. Thus the two gaps here
are one and two, and we can choose generators with
initial orders three, four, and five. They form an
ambient coordinate system: their classes generate the
three-dimensional cotangent space, since the square
of the maximal ideal has order at least six.
In a defining relation, the quadratic monomials
$x^2$ and $xy$ have respectively the unique least
orders six and seven, so their coefficients vanish.
Every gradient with valuation three consequently has
its initial vector in the single $z$ direction
(from a possible $xz$ term). Two such initial
gradient vectors are dependent. Every two-by-two
minor therefore has valuation at least seven,
contradicting $\tau_p=6$. Hence $d_p\ge7$ in this case.

\item Finally consider a nonplanar point with two branches
of multiplicities $m_1,m_2$. Two smooth curve branches
in a smooth threefold are contained in a smooth
surface. Indeed choose a common immersive coordinate
$x$ and write them as graphs
$(y,z)=(a_i(x),b_i(x))$. Divide their difference pair
by its common power of $x$; the equation
\[
 \frac{b_1-b_2}{x^k}(y-a_1)
 -\frac{a_1-a_2}{x^k}(z-b_1)=0
\]
has a nonzero linear derivative and contains both.
Thus in the nonplanar case $m_1+m_2\ge3$.
The absence of a linear defining equation gives,
on the two branches separately,
$\tau_p\ge2(m_1+m_2)$.
The reduced-curve normalization sequence gives
$\delta(C,p)=\delta_1+\delta_2+
\operatorname{length}\mathcal O_Y/(P_1+P_2)
\ge m_1+m_2-1$.
Also $r_p=m_1+m_2-2$.
It follows that
$d_p\ge2(m_1+m_2)+1\ge7$.
This proves the gap in all cases.

\item At the unique two-branch point, a planar singularity
has defect five if its delta invariant is one, and
at least ten otherwise. A nonplanar point has defect
at least seven. Every other singular point is
unibranch and contributes at least five if planar,
and at least seven if nonplanar. Smooth points
contribute zero. Thus the sum of the defects can
be five or at least seven, but never six.
The required value $\chi(\mathcal O_{C_{\rm sch}})=-1$
would give exactly six. This contradiction excludes
$e(S)=1$.

\item The final statement uses P4.26: the sum of all the
nonnegative fibre Euler numbers is at most two.
The rational-resolution case is strictly positive,
and we have just excluded one, leaving exactly two.

\end{enumerate}
\end{proof}
